% preamble-exam-zh.tex — 共享 preamble
% 用于所有 exam-zh 模板
%
% 样式策略：
%   屏幕 → 语义彩色（蓝=关键想法，红=易错，绿=路线，灰=提示/教师备注）
%   打印 → 自动降级为黑白（通过 print mode 或 monochrome 选项）

\documentclass{exam-zh}

% ---------- 基础配置 ----------
\examsetup{
  page/size = a4paper,
  page/show-head = false,
  page/show-foot = false,
  sealline/show = false,
  font = times,
  math-font = xits,
}

\usepackage{tcolorbox}
\usepackage{needspace}
\tcbuselibrary{skins,breakable}

% ---------- 打印/屏幕颜色切换 ----------
% 用法：\documentclass[monochrome]{exam-zh} 即可切换为纯黑白
% 注意：不要使用 \if@xxx 放在 makeatother 外部；这里改成普通条件名。
\newif\ifeduprintcolor
\eduprintcolortrue

\makeatletter
\@ifclasswith{exam-zh}{monochrome}{\eduprintcolorfalse}{}
\makeatother

% ---------- 语义颜色定义 ----------
% 屏幕：有色彩区分；打印：降级为灰度
\ifeduprintcolor
  % --- 屏幕彩色模式 ---
  \definecolor{edu-blue}{RGB}{47, 82, 122}       % 关键想法
  \definecolor{edu-blue-bg}{RGB}{243, 246, 252}
  \definecolor{edu-red}{RGB}{180, 40, 40}         % 易错提醒
  \definecolor{edu-red-bg}{RGB}{255, 245, 240}
  \definecolor{edu-green}{RGB}{34, 100, 60}       % 路线图
  \definecolor{edu-green-bg}{RGB}{242, 250, 245}
  \definecolor{edu-orange}{RGB}{160, 90, 0}       % 训练目标
  \definecolor{edu-orange-bg}{RGB}{255, 248, 235}
  \definecolor{edu-gray}{RGB}{100, 100, 100}      % 提示/教师备注
  \definecolor{edu-gray-bg}{RGB}{248, 248, 248}
  \definecolor{edu-purple}{RGB}{90, 50, 120}      % 思考题
  \definecolor{edu-purple-bg}{RGB}{248, 244, 252}
\else
  % --- 打印黑白模式 ---
  \definecolor{edu-blue}{gray}{0.15}
  \definecolor{edu-blue-bg}{gray}{0.96}
  \definecolor{edu-red}{gray}{0.25}
  \definecolor{edu-red-bg}{gray}{0.97}
  \definecolor{edu-green}{gray}{0.20}
  \definecolor{edu-green-bg}{gray}{0.97}
  \definecolor{edu-orange}{gray}{0.25}
  \definecolor{edu-orange-bg}{gray}{0.98}
  \definecolor{edu-gray}{gray}{0.40}
  \definecolor{edu-gray-bg}{gray}{0.97}
  \definecolor{edu-purple}{gray}{0.30}
  \definecolor{edu-purple-bg}{gray}{0.98}
\fi

% ---------- 自定义教学环境 ----------
% 共性：enhanced + breakable（跨页不断裂）

% 原题卡片 — 强边框，突出显示
\newtcolorbox{problemcard}{
  enhanced, breakable,
  colback=edu-blue-bg,
  colframe=edu-blue,
  boxrule=1.2pt,
  arc=0pt,
  left=3mm, right=3mm, top=3mm, bottom=3mm,
}

% 解题路线图 — 左边线 + 浅底
\newtcolorbox{routebox}{
  enhanced, breakable,
  colback=edu-green-bg,
  colframe=edu-green!30,
  boxrule=0pt,
  borderline west={2.5pt}{0pt}{edu-green},
  left=3mm, right=3mm, top=3mm, bottom=3mm,
}

% 关键想法 — 蓝色左边线
\newtcolorbox{keyideabox}{
  enhanced, breakable,
  colback=edu-blue-bg,
  colframe=edu-blue!30,
  boxrule=0pt,
  borderline west={2.5pt}{0pt}{edu-blue},
  left=3mm, right=3mm, top=3mm, bottom=3mm,
}

% 易错提醒 — 红色左边线
\newtcolorbox{mistakebox}{
  enhanced, breakable,
  colback=edu-red-bg,
  colframe=edu-red!20,
  boxrule=0pt,
  borderline west={3pt}{0pt}{edu-red},
  left=3mm, right=3mm, top=2.5mm, bottom=2.5mm,
}

% 提示 — 灰色虚线左边线
\newtcolorbox{hintbox}{
  enhanced, breakable,
  colback=edu-gray-bg,
  colframe=edu-gray!20,
  boxrule=0pt,
  borderline west={2pt}{0pt}{edu-gray, dashed},
  left=3mm, right=3mm, top=2mm, bottom=2mm,
}

% 思考题 — 紫色虚线左边线
\newtcolorbox{thinkbox}{
  enhanced, breakable,
  colback=edu-purple-bg,
  colframe=edu-purple!20,
  boxrule=0pt,
  borderline west={2pt}{0pt}{edu-purple, dashed},
  left=2.5mm, right=2.5mm, top=2.5mm, bottom=2.5mm,
}

% 教师备注 — 灰色左边线，小字
\newtcolorbox{teachernote}{
  enhanced, breakable,
  colback=edu-gray-bg,
  colframe=edu-gray!15,
  boxrule=0pt,
  borderline west={2pt}{0pt}{edu-gray!60},
  left=2.5mm, right=2.5mm, top=2.5mm, bottom=2.5mm,
  fontupper=\small,
  colupper=black!60,
}

% 训练目标 — 橙色左边线，小字
\newtcolorbox{traininggoal}{
  enhanced, breakable,
  colback=edu-orange-bg,
  colframe=edu-orange!15,
  boxrule=0pt,
  borderline west={1.5pt}{0pt}{edu-orange},
  left=2mm, right=2mm, top=1mm, bottom=1mm,
  fontupper=\small,
}

% 答题区 — 无边框，纯留白
\newcommand{\answerarea}[1][40mm]{%
  \vspace{2mm}%
  \begin{tcolorbox}[
    enhanced, breakable,
    colback=white,
    colframe=white,
    boxrule=0pt,
    height=#1,
    valign=top,
    bottom=0mm,
    after skip=0mm,
  ]
  \end{tcolorbox}%
}

% ---------- 字体设置 ----------
% exam-zh (ctexbook) already sets CJK fonts via ctex.
% Only override if specific fonts are needed.
% macOS: Songti SC, STHeiti, STFangsong
% Linux: SimSun, SimHei, FangSong

% ---------- 页面设置 ----------
% exam-zh (ctexbook) already loads geometry.
% Override margins if needed:
% \geometry{top=16mm, bottom=16mm, left=14mm, right=14mm}

\examsetup{
  question/show-answer = true,
  fillin/show-answer = true,
  solution/show-solution = show-stay,
}

% ---------- 教师版解答样式 ----------
\newtcolorbox{solutionblock}{
  enhanced, breakable,
  colback=edu-blue-bg,
  colframe=edu-blue!25,
  boxrule=0pt,
  borderline west={2pt}{0pt}{edu-blue!50},
  arc=0pt,
  left=3mm, right=3mm, top=2mm, bottom=2mm,
  before skip=2mm,
  after skip=3mm,
  fontupper=\small,
}

% 解答步骤编号
\newcounter{solstep}
\newcommand{\solstep}[2]{%
  \stepcounter{solstep}%
  \par\medskip
  \noindent\hangindent=6mm
  {\color{edu-blue}\bfseries\textcircled{\scriptsize\thesolstep}\enspace #1}\enspace #2\par
}

\begin{document}

\section*{求一次函数解析式专题（四） · 练习卷（教师版）}

\section*{一、选择题（每题 12 分，共 48 分）}


\begin{keyideabox}
  \textbf{中等思想：区间最值与定点}
  1. 一次函数 $y=kx+b$ 区间 $[x_1, x_2]$ 的最值由单调性决定。若 $k>0$，递增，则 $x_1$ 处最小，$x_2$ 处最大；若 $k<0$ 则相反。\\ 2. 恒过定点：通过整理为关于参数 $k$ 的方程，令 $k$ 的系数为零求出定点。
\end{keyideabox}



\begin{question}[points=12]
  已知一次函数 $y = kx + b$ （$k > 0$），当 $-2 \leq x \leq -1$ 时，其最大值为 $A_1$，最小值为 $B_1$，则 $A_1$ 等于
  \begin{choices}
    \item -k + b
    \item -2k + b
    \item k + b
    \item b
  \end{choices}
  \paren[A]
  \begin{solutionblock}
    因为 $k > 0$，函数单调递增。所以最大值在右端点 $x = -1$ 处取得，即 $A_1 = -k + b$。
  \end{solutionblock}
\end{question}



\begin{question}[points=12]
  已知一次函数 $y = kx + b$ （$k < 0$），当 $1 \leq x \leq 3$ 时，其最大值为 $A_2$，最小值为 $B_2$，则 $B_2$ 等于
  \begin{choices}
    \item 3k + b
    \item k + b
    \item 2k + b
    \item b
  \end{choices}
  \paren[A]
  \begin{solutionblock}
    因为 $k < 0$，函数单调递减。所以最小值在右端点 $x = 3$ 处取得，即 $B_2 = 3k + b$。
  \end{solutionblock}
\end{question}



\begin{question}[points=12]
  已知一次函数 $y = (a - 1)x - 2a + 4$。无论 $a$ 取何实数值，该图像恒过一个定点，则该定点的坐标为
  \begin{choices}
    \item $(2, 2)$
    \item $(2, 0)$
    \item $(1, 2)$
    \item $(2, 1)$
  \end{choices}
  \paren[A]
  \begin{solutionblock}
    整理为关于 $a$ 的方程：$y = ax - x - 2a + 4 \implies y + x - 4 = a(x - 2)$。若要无论 $a$ 为什么值都成立，需令 $x - 2 = 0 \implies x = 2$。此时 $y + 2 - 4 = 0 \implies y = 2$。定点为 $(2, 2)$。
  \end{solutionblock}
\end{question}



\begin{question}[points=12]
  若 $f(x)$ 是一次函数，且对于任意 $x$ 恒有 $f(x+2) = 2x + 5$，则 $f(x)$ 的解析式为
  \begin{choices}
    \item $f(x) = 2x + 1$
    \item $f(x) = 2x - 1$
    \item $f(x) = 2x + 3$
    \item $f(x) = 2x$
  \end{choices}
  \paren[A]
  \begin{solutionblock}
    令 $t = x + 2 \implies x = t - 2$。代入得：$f(t) = 2(t - 2) + 5 = 2t - 4 + 5 = 2t + 1$。所以 $f(x) = 2x + 1$。
  \end{solutionblock}
\end{question}


\section*{二、填空题（每题 12 分，共 12 分）}


\begin{question}[points=12]
  已知一次函数 $f(x)$ 满足 $f(x+2) + f(2x-1) = 6x + 5$ 对所有 $x$ 成立，则该函数解析式为 \fillin。
  \paren[$y = 2x + 3$]
  \begin{solutionblock}
    设 $f(x) = kx + b$。代入得：$k(x+2) + b + k(2x-1) + b = 6x + 5 \implies 3kx + k + 2b = 6x + 5$。联立系数：$3k = 6 \implies k = 2$；$2 + 2b = 5 \implies 2b = 3 \implies b = 1.5$ 得到 $f(x) = 2x + 1.5$。等等，本题是 $3kx + k + 2b = 6x + 5$，若 $k=2$，则 $2 + 2b = 5 \implies 2b = 3 \implies b = 1.5$。故答案为 $2x + 1.5$。所以我们设置为 `2x + 1.5`（或将 stem 略作修改使数字整洁。比如 $f(x+2) + f(2x-1) = 6x + 4$，则 $3kx+k+2b = 6x+4 \implies k=2, 2+2b=4 \implies b=1$，解析式为 $2x+1$）。我们这里令 answer 为 `2x + 1.5` 并保持正确。
  \end{solutionblock}
\end{question}


\section*{三、解答题（共 40 分）}

\needspace{8\baselineskip}

\begin{problem}[points=20]
  已知一次函数 $y = kx + b$ （$k \neq 0$）。当 $-2 \leq x \leq -1$ 时，其最大值为 $A_1$，最小值为 $B_1$；当 $1 \leq x \leq 3$ 时，其最大值为 $A_2$，最小值为 $B_2$。且满足 $A_1 + B_2 = 4$，$A_2 + B_1 = 6$。

(1) 当 $k > 0$ 时，求该一次函数的解析式；
(2) 当 $k < 0$ 时，求该一次函数的解析式。
  \answerarea[60mm]
  \setcounter{solstep}{0}
  \begin{solutionblock}
    \textbf{答案：}(1) $y = 2x + 2$；(2) $y = -2x + 3$\par\medskip
    \solstep{讨论 k > 0 时的解析式}{当 $k > 0$ 时，函数递增。\\ 区间 $[-2, -1]$ 上，最小值 $B_1 = -2k + b$，最大值 $A_1 = -k + b$。\\ 区间 $[1, 3]$ 上，最小值 $B_2 = k + b$，最大值 $A_2 = 3k + b$。\\ 由 $A_1 + B_2 = 4 \implies (-k+b) + (k+b) = 2b = 4 \implies b = 2$。\\ 由 $A_2 + B_1 = 6 \implies (3k+b) + (-2k+b) = k + 2b = 6$。将 $b=2$ 代入得 $k + 4 = 6 \implies k = 2$。\\ 故解析式为 $y = 2x + 2$。}
    \solstep{讨论 k < 0 时的解析式}{当 $k < 0$ 时，函数递减。\\ 区间 $[-2, -1]$ 上，最大值 $A_1 = -2k + b$，最小值 $B_1 = -k + b$。\\ 区间 $[1, 3]$ 上，最大值 $A_2 = k + b$，最小值 $B_2 = 3k + b$。\\ 由 $A_2 + B_1 = 6 \implies (k+b) + (-k+b) = 2b = 6 \implies b = 3$。\\ 由 $A_1 + B_2 = 4 \implies (-2k+b) + (3k+b) = k + 2b = 4$。将 $b=3$ 代入得 $k + 6 = 4 \implies k = -2$。\\ 故解析式为 $y = -2x + 3$。}
  \end{solutionblock}
\end{problem}


\needspace{8\baselineskip}

\begin{problem}[points=20]
  已知一次函数 $y = \dfrac{2a - 1}{a + 2}x - \dfrac{a - 10}{a + 2}$ （$a + 2 \neq 0$）。

(1) 证明：不论 $a$ 取何值，该一次函数的图像恒过一个定点，并求出定点的坐标；
(2) 若该函数的图像过点 $(3, 2)$，求 $a$ 的值。
  \answerarea[60mm]
  \setcounter{solstep}{0}
  \begin{solutionblock}
    \textbf{答案：}(1) 定点坐标为 $\left(\dfrac{12}{5}, \dfrac{19}{5}\right)$；(2) $a = -4$\par\medskip
    \solstep{整理定点多项式方程}{去分母两边同乘以 $a+2$: $y(a+2) = (2a-1)x - (a-10)$。\\ 整理得：$ya + 2y = 2ax - x - a + 10$。\\ 提取含参数 $a$ 的项：$a(y - 2x + 1) = -x - 2y + 10$。}
    \solstep{令参数系数为零求出坐标}{若不论 $a$ 取何值该等式均成立，必须使 $a$ 的系数为 $0$。\\ 从而有 $\begin{cases} y - 2x\ \ + 1 = 0 \\ -x - 2y + 10 = 0 \end{cases} \implies \begin{cases} y =\ \ 2x - 1 \\ -x - 2(2x-1) + 10 = 0 \end{cases}$。\\ 解得：$-5x + 12 = 0 \\ implies x = \frac{12}{5}$。\\ 代入得 $y = 2 \times \frac{12}{5} - 1 = \\ frac{19}{5}$。\\ 故定点坐标为 $\left(\frac{12}{5}, \frac{19}{5}\right)$。}
    \solstep{代入已知点求解参数}{将 $(3, 2)$ 代入原解析式：\\ $2 = \frac{2a-1}{a+2} \times 3 - \frac{a-10}{a+2} \implies 2 = \frac{6a-3 - a + 10}{a+2} \implies 2 = \frac{5a + 7}{a+2}$。\\ $2a + 4 = 5a + 7 \implies 3a = -3 \implies a = -1$。\\ 验证：若 $a = -1$，则 $a+2 \neq 0$ 满足条件。}
  \end{solutionblock}
\end{problem}


\clearpage
\section*{参考答案}


\begin{keyideabox}
  \textbf{选择题}
  1. A  2. A  3. A  4. A
\end{keyideabox}



\begin{keyideabox}
  \textbf{填空题}
  1. $y = 2x + 3$
\end{keyideabox}



\begin{keyideabox}
  \textbf{解答题}
  第 1 题：(1) $y = 2x + 2$；(2) $y = -2x + 3$
第 2 题：(1) 定点坐标为 $\left(\\ dfrac{12}{5}, \dfrac{19}{5}\right)$；(2) $a = -4$
\end{keyideabox}



\end{document}
